\begin{theorem}
$\sqrt{7}$ is irrational.
\end{theorem}

\begin{proof}
Suppose, for contradiction, that $\sqrt{7}$ is rational. Then there exist integers $a$ and $b$ with $b \neq 0$ such that
\[
\sqrt{7}=\frac{a}{b},
\]
and we may choose $a$ and $b$ so that the fraction is in lowest terms, meaning
\[
\gcd(a,b)=1.
\]
(Such a representation exists because, given any integers $a$ and $b$ with $b \neq 0$, we can divide both by $\gcd(a,b)$ to obtain an equivalent fraction whose numerator and denominator are coprime.)

Squaring both sides gives
\[
7=\frac{a^2}{b^2}.
\]

Multiplying by $b^2$, we get
\[
a^2=7b^2.
\]

Thus $7 \mid a^2$. Since $7$ is prime, if $7 \mid a^2$, then $7 \mid a$ (by Euclid's lemma: if a prime $p$ divides a product $mn$, then $p \mid m$ or $p \mid n$; here applied with $m = n = a$). Therefore, there exists an integer $k$ such that
\[
a=7k.
\]

Substituting this into $a^2=7b^2$, we obtain
\[
(7k)^2=7b^2,
\]
hence
\[
49k^2=7b^2.
\]

Dividing by $7$ gives
\[
7k^2=b^2.
\]

So $7 \mid b^2$. Again, since $7$ is prime, Euclid's lemma implies $7 \mid b$.

Therefore, both $a$ and $b$ are divisible by $7$, so $\gcd(a,b) \geq 7$. This contradicts the assumption that $\gcd(a,b)=1$.

This contradiction shows that $\sqrt{7}$ cannot be rational; hence $\sqrt{7}$ is irrational.
\end{proof}
